\documentclass[11pt,openany]{report}
\usepackage[letterpaper,margin=0.82in,headheight=14pt]{geometry}
\usepackage[T1]{fontenc}
\usepackage[utf8]{inputenc}
\usepackage{lmodern}
\usepackage{microtype}
\usepackage{amsmath,amssymb,amsthm,mathtools}
\usepackage{booktabs,longtable,tabularx,array,multirow}
\usepackage{graphicx,float,placeins}
\usepackage[dvipsnames]{xcolor}
\usepackage{enumitem}
\usepackage{fancyhdr}
\usepackage{listings}
\usepackage{url}
\usepackage[hidelinks]{hyperref}
\usepackage[nameinlink,noabbrev]{cleveref}
\usepackage{needspace}
\allowdisplaybreaks

\hypersetup{
  pdftitle={Projector-Generated Special Functions VIII},
  pdfauthor={Komorebi P.},
  pdfsubject={q-Hahn lowest-weight packets, q-Racah recoupling, Askey-Wilson transfer functions, and deformation geometry}
}

\definecolor{Ink}{HTML}{17202A}
\definecolor{Blue}{HTML}{2E86AB}
\definecolor{Teal}{HTML}{2A9D8F}
\definecolor{Orange}{HTML}{F4A261}
\definecolor{Magenta}{HTML}{A23E8C}
\definecolor{Red}{HTML}{E76F51}
\definecolor{Paper}{HTML}{FAFAF7}
\definecolor{Muted}{HTML}{5D6D7E}
\definecolor{CodeBG}{HTML}{F3F5F5}

\pagecolor{Paper}
\color{Ink}
\graphicspath{{figures/}}
\setlength{\parindent}{0pt}
\setlength{\parskip}{0.52em}
\raggedbottom

\pagestyle{fancy}
\fancyhf{}
\fancyhead[L]{\small Projector-Generated Special Functions VIII}
\fancyhead[R]{\small Komorebi P.}
\fancyfoot[C]{\thepage}
\renewcommand{\headrulewidth}{0.4pt}

\newtheorem{theorem}{Theorem}[chapter]
\newtheorem{proposition}[theorem]{Proposition}
\newtheorem{lemma}[theorem]{Lemma}
\newtheorem{corollary}[theorem]{Corollary}
\newtheorem{conjecture}[theorem]{Conjecture}
\theoremstyle{definition}
\newtheorem{definition}[theorem]{Definition}
\newtheorem{example}[theorem]{Worked example}
\newtheorem{algorithm}[theorem]{Algorithm}
\theoremstyle{remark}
\newtheorem{remark}[theorem]{Remark}
\newtheorem{boundary}[theorem]{Boundary statement}

\newcommand{\C}{\mathbb C}
\newcommand{\R}{\mathbb R}
\newcommand{\N}{\mathbb N}
\newcommand{\e}{\mathrm e}
\newcommand{\ii}{\mathrm i}
\newcommand{\dd}{\,\mathrm d}
\newcommand{\norm}[1]{\left\lVert#1\right\rVert}
\newcommand{\abs}[1]{\left\lvert#1\right\rvert}
\newcommand{\diag}{\operatorname{diag}}
\newcommand{\Id}{\mathrm I}
\newcommand{\Tr}{\operatorname{Tr}}
\newcommand{\rank}{\operatorname{rank}}
\newcommand{\spanop}{\operatorname{span}}
\newcommand{\Cas}{\mathcal C}
\newcommand{\Mcal}{\mathcal M}
\newcommand{\Rcal}{\mathcal R}
\newcommand{\Pcal}{\mathcal P}
\newcommand{\Hcal}{\mathcal H}
\newcommand{\ket}[1]{\left|#1\right\rangle}
\newcommand{\bra}[1]{\left\langle#1\right|}
\newcommand{\ip}[2]{\left\langle#1,#2\right\rangle}
\newcommand{\qnum}[1]{\left[#1\right]_q}
\newcommand{\qpoch}[2]{\left(#1;q\right)_{#2}}
\newcommand{\qrise}[2]{\left[#1\right]_{q,#2}}
\newcommand{\qbinom}[2]{\genfrac{[}{]}{0pt}{}{#1}{#2}_{q}}
\newcommand{\entropy}{\operatorname{H}}
\newcommand{\Res}{\operatorname{Res}}

\newcolumntype{L}[1]{>{\raggedright\arraybackslash}p{#1}}

\lstdefinestyle{pythonstyle}{
  language=Python,
  basicstyle=\ttfamily\scriptsize,
  backgroundcolor=\color{CodeBG},
  frame=single,
  rulecolor=\color{Muted!45},
  breaklines=true,
  breakatwhitespace=false,
  showstringspaces=false,
  numbers=left,
  numberstyle=\tiny\color{Muted},
  numbersep=7pt,
  tabsize=4,
  columns=fullflexible,
  keepspaces=true
}
\lstset{style=pythonstyle}

\newcommand{\FigureWide}[3]{%
\Needspace{0.52\textheight}
\begin{figure}[H]
\centering
\includegraphics[width=0.92\textwidth,height=0.43\textheight,keepaspectratio]{#1.pdf}
\caption{#2 #3}
\label{fig:#1}
\end{figure}
}

\newcommand{\InlinePart}[1]{%
\FloatBarrier
\Needspace{10\baselineskip}
\refstepcounter{part}%
\addcontentsline{toc}{part}{\protect\numberline{\thepart}#1}%
\vspace{1.4em}
{\color{Blue}\Large\bfseries Part \thepart\par}
\vspace{0.2em}
{\color{Blue}\LARGE\bfseries #1\par}
\vspace{0.9em}
}

\newcommand{\FlowChapter}[1]{%
\FloatBarrier
\Needspace{14\baselineskip}
\refstepcounter{chapter}%
\setcounter{section}{0}%
\setcounter{subsection}{0}%
\addcontentsline{toc}{chapter}{\protect\numberline{\thechapter}#1}%
\markboth{#1}{#1}%
\vspace{1.3em}
{\color{Muted}\large\bfseries Chapter \thechapter\par}
\vspace{0.15em}
{\LARGE\bfseries #1\par}
\vspace{0.9em}
}

\title{\textbf{Projector-Generated Special Functions VIII}\\[0.25em]
\Large q-Hahn Lowest-Weight Packets, q-Racah Recoupling, Askey-Wilson Transfer Functions, and Deformation Geometry}
\author{Komorebi P.}
\date{}

\begin{document}
\hypersetup{pageanchor=false}
\pagenumbering{roman}
\maketitle

\begin{abstract}
Papers VI and VII produced a continuous shape semigroup, positive discrete-series packet bases, classical Racah recoupling matrices, bistochastic measurement kernels, and pentagon coherence. The present paper introduces a positive deformation parameter $q$ and rebuilds the same architecture inside a symmetric $U_q(\mathfrak{su}(1,1))$ coordinate system.

For $0<q\le1$, define the symmetric q-number
\[
\qnum{x}=\frac{q^x-q^{-x}}{q-q^{-1}},
\]
with the continuous value $\qnum{x}=x$ at $q=1$. The positive shape packet $D^+_{\kappa,q}$ carries the ladder
\[
K_0\ket n=(n+\kappa)\ket n,
\qquad
K_+\ket n=\sqrt{\qnum{n+1}\qnum{2\kappa+n}}\ket{n+1},
\]
\[
K_-\ket n=\sqrt{\qnum n\qnum{2\kappa+n-1}}\ket{n-1}.
\]
The coproduct
\[
\Delta K_-=K_-\otimes q^{K_0}+q^{-K_0}\otimes K_-
\]
produces exact fixed-excitation packet spaces and a q-Casimir. The lowest-weight vector of two packets has an explicit q-Hahn coefficient formula. Three packets generate two intermediate q-Casimir atlases whose overlap is orthogonal and belongs to the q-Racah class. Squaring the overlap produces a doubly stochastic measurement kernel.

For three equal half-shape packets at one excitation, the complete transform is
\[
\boxed{
\Rcal_q=
\begin{pmatrix}
\dfrac{q}{1+q^2}&\dfrac{\sqrt{1+q^2+q^4}}{1+q^2}\\[0.8em]
\dfrac{\sqrt{1+q^2+q^4}}{1+q^2}&-\dfrac{q}{1+q^2}
\end{pmatrix}.
}
\]
It converges to the classical Racah beam splitter at $q=1$ and to a routing swap as $q\downarrow0$. The corresponding row entropy and mixing gap are monotone in this two-channel case. A higher-excitation counterexample shows that such monotonicity is not universal.

Writing $q=\e^{-\theta}$ turns the recoupling matrices into a differentiable orthogonal flow. In a continuous sign gauge, the generator
\[
A_N(\theta)=\Rcal_N(\theta)^{\mathsf T}\partial_\theta\Rcal_N(\theta)
\]
is skew-symmetric. Its Frobenius norm defines a finite deformation metric. For four packets, five q-Casimir tree bases form the associahedral pentagon, and their ordered overlap product is the identity. Every finite q-Racah packet matrix can be compiled into Givens rotations and diagonal signs.

The scalar side is developed independently through the terminating q-Racah ${}_4\phi_3$ function and the continuous Askey-Wilson function. Exact duality, the $q\to1$ Racah limit, and the Askey-Wilson q-difference equation are verified symbolically or at high precision. q-Racah and Askey-Wilson functions, quantum-group recoupling, and the pentagon identity are classical. The generated contribution is the common packet architecture that retains the q-Hahn seed, projector reduction, q-Racah analysis and synthesis, stochastic measurement, deformation flow, associator certificate, and finite compiler in one system.

All figures were generated by the accompanying Python program. The release contains 32 vector figures, matching PNG files, 23 built-in checks, 20 independent tests, numerical tables, and the complete source.
\end{abstract}

\chapter*{Status and claim boundary}
\addcontentsline{toc}{chapter}{Status and claim boundary}
The q-Racah and Askey-Wilson polynomial families are classical members of the Askey-Wilson scheme. Their q-hypergeometric definitions, orthogonality measures, duality laws, q-difference equations, and classical limits are established parts of special-function theory. Their occurrence as Clebsch-Gordan or Racah coefficients in quantum-group representation theory is also classical.

This paper does not claim to have invented q-Racah polynomials, Askey-Wilson polynomials, quantum $6j$ symbols, or the abstract pentagon identity. It shows how they arise from the projector-resolved packet construction and adds a single operational language for lowest-weight reduction, analysis, synthesis, measurement, deformation transport, and compilation.

The finite matrix construction uses a symmetric q-number convention and the positive real range $0<q\le1$. Other conventions are related by changes of generators and powers of $q$. Formula comparisons must therefore preserve the declared convention.

The q-dependent entropy statements concern the classical channel obtained after squaring the amplitudes and discarding phase. The underlying q-Racah matrix remains exactly reversible. Entropy is monotone under each fixed bistochastic channel. It is not generally monotone as the deformation parameter itself changes.

The four-body pentagon closes because every tree basis is complete in the same lowest-weight multiplicity space. The numerical product is a stringent implementation certificate, not a replacement for the abstract Biedenharn-Elliott theorem.

The final quantum discussion keeps this shape-recoupling module separate from the logarithmic number Hamiltonian used in the attached zeta-governed quantum dynamics. The modules can be tensored or composed, but the present paper does not identify their spectra and does not prove the Riemann hypothesis.

\begingroup
\footnotesize
\setlength{\parskip}{0pt}
\tableofcontents
\listoffigures
\endgroup
\clearpage
\pagenumbering{arabic}
\hypersetup{pageanchor=true}

\InlinePart{q-deformed packet algebra}

\FlowChapter{Symmetric q-numbers and positive shape packets}

\begin{definition}[Symmetric q-number]
For $q>0$ and $q\ne1$, define
\[
\boxed{
\qnum{x}=\frac{q^x-q^{-x}}{q-q^{-1}}.
}
\]
At $q=1$, define $\qnum{x}=x$ by continuity.
\end{definition}

Writing $q=\e^{-\theta}$ gives
\[
\qnum{x}=\frac{\sinh(x\theta)}{\sinh\theta}.
\]
Hence $\qnum{-x}=-\qnum{x}$ and $\qnum{x}>0$ for $x>0$.

\begin{definition}[Positive q-shape packet]
For $\kappa>0$, let $D^+_{\kappa,q}$ have orthonormal basis $\{\ket n:n\ge0\}$ with
\[
K_0\ket n=(n+\kappa)\ket n,
\]
\[
K_+\ket n=
\sqrt{\qnum{n+1}\qnum{2\kappa+n}}\ket{n+1},
\]
\[
K_-\ket n=
\sqrt{\qnum n\qnum{2\kappa+n-1}}\ket{n-1}.
\]
\end{definition}

\begin{proposition}[Local q-algebra]
On every basis vector away from a finite truncation boundary,
\[
[K_0,K_\pm]=\pm K_\pm,
\qquad
[K_-,K_+]=\qnum{2K_0}.
\]
\end{proposition}

\begin{proof}
The first identities follow because $K_\pm$ change $n$ by one. For the second, the coefficient of $\ket n$ in $K_-K_+\ket n$ is
\[
\qnum{n+1}\qnum{2\kappa+n},
\]
while the coefficient in $K_+K_-\ket n$ is
\[
\qnum n\qnum{2\kappa+n-1}.
\]
The elementary hyperbolic identity
\[
\qnum{x+1}\qnum y-\qnum x\qnum{y-1}=\qnum{x+y}
\]
with $x=n$ and $y=2\kappa+n$ gives $\qnum{2n+2\kappa}=\qnum{2K_0}$.
\end{proof}

\FigureWide{01_theory_pipeline}{The q-deformed packet program.}{The positive-series ladder generates a q-Hahn lowest packet, pair q-Casimir atlases, q-Racah transforms, measurement kernels, deformation flow, and four-body associator closure.}

\FigureWide{02_q_number_deformation}{Symmetric q-numbers.}{As $q$ approaches one, the deformed ladder approaches the ordinary linear coordinate. Smaller $q$ amplifies large positive degrees.}

\FigureWide{03_q_ladder_coefficients}{Positive-series raising coefficients.}{The local q-ladder remains exact while its high-degree growth changes strongly with the deformation parameter.}

\FlowChapter{Coproduct, q-Casimir, and exact tensor routing}

Define
\[
\Delta K_0=K_0\otimes\Id+\Id\otimes K_0,
\]
\[
\boxed{
\Delta K_\pm
=K_\pm\otimes q^{K_0}+q^{-K_0}\otimes K_\pm.
}
\]
Iterating this formula gives the total generators on any finite packet product.

\begin{theorem}[Coassociative packet routing]
The coproduct is coassociative on the declared generators. In particular,
\[
(\Delta\otimes\Id)\Delta K_-
=(\Id\otimes\Delta)\Delta K_-
\]
and both sides equal
\[
K_-\otimes q^{K_0}\otimes q^{K_0}
+q^{-K_0}\otimes K_-\otimes q^{K_0}
+q^{-K_0}\otimes q^{-K_0}\otimes K_-.
\]
\end{theorem}

\begin{proof}
Since $K_0$ is primitive,
\[
q^{\Delta K_0}=q^{K_0}\otimes q^{K_0}.
\]
Substituting this identity into either iterated coproduct gives the same three terms. The proof for $K_+$ is identical.
\end{proof}

\begin{definition}[q-Casimir]
Define
\[
\boxed{
\Cas_q=\qnum{K_0}\qnum{K_0-1}-K_+K_-.
}
\]
\end{definition}

\begin{theorem}[Scalar packet Casimir]
On $D^+_{\kappa,q}$,
\[
\boxed{
\Cas_q=\qnum\kappa\qnum{\kappa-1}\Id.
}
\]
\end{theorem}

\begin{proof}
On $\ket n$ the first term is
\[
\qnum{n+\kappa}\qnum{n+\kappa-1},
\]
and the second is
\[
\qnum n\qnum{2\kappa+n-1}.
\]
The q-number identity
\[
\qnum{x+y}\qnum{x+y-1}-\qnum y\qnum{2x+y-1}
=\qnum x\qnum{x-1}
\]
with $x=\kappa$ and $y=n$ gives the result.
\end{proof}

\FigureWide{04_q_coproduct_paths}{Coproduct paths.}{A single lowering event is routed through each factor with the q-weights required by coassociativity.}

\FlowChapter{Fixed excitation and the lowest-weight multiplicity packet}

For $m$ factors, let
\[
\Hcal_N^{(m)}
=\spanop\{\ket{n_1,\ldots,n_m}:n_1+\cdots+n_m=N\}.
\]
Then
\[
\dim\Hcal_N^{(m)}=\binom{N+m-1}{m-1}.
\]
The lowest-weight multiplicity packet is
\[
\boxed{
\Mcal_N^{(m)}
=\ker\left(K_-^{[m]}\big|_{\Hcal_N^{(m)}}\right).
}
\]

\begin{theorem}[Multiplicity dimension]
For every $N\ge0$ and $m\ge2$,
\[
\boxed{
\dim\Mcal_N^{(m)}=\binom{N+m-2}{m-2}.
}
\]
\end{theorem}

\begin{proof}
The lowering map sends $\Hcal_N^{(m)}$ to $\Hcal_{N-1}^{(m)}$. Order the target compositions by their last coordinate. For every target state $t$, choose the source $t+e_m$. Lowering the last factor produces $t$ with a nonzero coefficient. Lowering any earlier factor produces a target with a larger last coordinate. The corresponding square submatrix is triangular with nonzero diagonal, so the lowering map is onto. Rank-nullity gives
\[
\dim\Mcal_N^{(m)}
=\binom{N+m-1}{m-1}-\binom{N+m-2}{m-1}
=\binom{N+m-2}{m-2}.
\]
\end{proof}

This projector reduction is output-sensitive. The full tensor level defines the kernel. All recoupling then occurs in the smaller multiplicity packet.

\InlinePart{q-Hahn lowest-weight functions}

\FlowChapter{Two-packet seed and closed coefficients}

Take two shapes $\kappa_1,\kappa_2>0$ and total excitation $N$. Write the fixed-excitation basis as
\[
\ket{n,N-n},
\qquad
0\le n\le N.
\]
Set
\[
\lambda=q^{-(N-1+\kappa_1+\kappa_2)}.
\]

\begin{theorem}[Explicit q-Hahn lowest packet]
A normalized lowest-weight vector is
\[
\boxed{
\ket{\Omega_N}
=\frac1{\sqrt{Z_N}}
\sum_{n=0}^{N}
(-\lambda)^n
\sqrt{
\qbinom{N}{n}
\frac{
\qrise{2\kappa_2+N-n}{n}
}{
\qrise{2\kappa_1}{n}
}
}
\ket{n,N-n},
}
\]
where
\[
\qbinom{N}{n}
=\frac{\qnum{N}!}{\qnum n!\qnum{N-n}!}
\]
and
\[
\qrise{x}{n}=\prod_{j=0}^{n-1}\qnum{x+j}.
\]
\end{theorem}

\begin{proof}
Let $c_n$ be the coefficient of $\ket{n,N-n}$. The coefficient of $\ket{n,N-n-1}$ in the total lowering equation is
\[
q^{N-n-1+\kappa_2}
\sqrt{\qnum{n+1}\qnum{2\kappa_1+n}}c_{n+1}
+
q^{-(n+\kappa_1)}
\sqrt{\qnum{N-n}\qnum{2\kappa_2+N-n-1}}c_n.
\]
Setting this to zero gives
\[
\frac{c_{n+1}}{c_n}
=-\lambda
\sqrt{
\frac{
\qnum{N-n}\qnum{2\kappa_2+N-n-1}
}{
\qnum{n+1}\qnum{2\kappa_1+n}
}
}.
\]
The displayed closed form has exactly this ratio. Normalization fixes $Z_N$.
\end{proof}

The squared coefficients form the discrete weight for the lowest q-Hahn packet. General coupled states are obtained by applying the total raising operator and are described by q-Hahn polynomial coordinates in the standard representation-theoretic normalization.

\FigureWide{05_q_hahn_packet_distribution}{q-Hahn lowest-weight packets.}{The same total excitation is distributed differently as q changes. The alternating phase is omitted in the probability plot.}

\FigureWide{06_q_hahn_classical_limit}{Classical limit of the lowest packet.}{The q-Hahn coefficient vector converges to its ordinary Hahn counterpart as q approaches one.}

\begin{example}[A five-excitation seed]
For
\[
N=5,
\qquad
(\kappa_1,\kappa_2)=(0.7,1.1),
\qquad
q=0.73,
\]
the code constructs the six coefficients from the theorem and verifies
\[
\norm{K_-^{[2]}\Omega_5}_\infty
=2.22\times10^{-16}.
\]
The numerical residual is not the proof. It confirms that the implementation follows the exact recurrence.
\end{example}

\InlinePart{q-Racah recoupling}

\FlowChapter{Two intermediate q-Casimir atlases}

For three shapes, the multiplicity space has dimension $N+1$. Restrict the pair Casimirs to this space:
\[
\Cas_{12}^{(N)}
=Q_N^{\mathsf T}\Cas_{12}Q_N,
\qquad
\Cas_{23}^{(N)}
=Q_N^{\mathsf T}\Cas_{23}Q_N,
\]
where the columns of $Q_N$ form an orthonormal frame for $\Mcal_N^{(3)}$.

The pair spectra are
\[
\boxed{
\lambda_j^{ab}(q)
=
\qnum{\kappa_a+\kappa_b+j}
\qnum{\kappa_a+\kappa_b+j-1},
\qquad
0\le j\le N.
}
\]
Let $L_N$ diagonalize $\Cas_{12}^{(N)}$ and let $R_N$ diagonalize $\Cas_{23}^{(N)}$. The recoupling matrix is
\[
\boxed{
\Rcal_N(q)=L_N^{\mathsf T}R_N.
}
\]

\begin{theorem}[Exact analysis and synthesis]
For every admissible finite packet,
\[
\boxed{
\Rcal_N(q)^{\mathsf T}\Rcal_N(q)
=
\Rcal_N(q)\Rcal_N(q)^{\mathsf T}
=
\Id_{N+1}.
}
\]
\end{theorem}

\begin{proof}
Both $L_N$ and $R_N$ are orthogonal eigenbases of real symmetric operators in the same multiplicity space. Therefore
\[
\Rcal_N^{\mathsf T}\Rcal_N
=R_N^{\mathsf T}L_NL_N^{\mathsf T}R_N
=\Id.
\]
The other identity is identical.
\end{proof}

The classical representation-theoretic identification says that, after normalization and a parameter map determined by the three shapes and $N$, the entries of $\Rcal_N$ are q-Racah functions. Equivalently, the restricted pair Casimirs form a finite tridiagonal pair. In the left atlas,
\[
\Rcal_N\diag(\lambda_m^{23})\Rcal_N^{\mathsf T}
\]
is tridiagonal. Hence each column obeys a three-term recurrence in the left packet label.

\FigureWide{10_q_casimir_spectra}{Pair q-Casimir spectra.}{The intermediate eigenvalues are exact q-number deformations of the classical quadratic spectrum.}

\FigureWide{11_q_racah_heatmaps}{q-Racah recoupling matrices.}{Changing q redistributes amplitude across the same complete pair-Casimir coordinate systems.}

\FigureWide{13_q_orthogonality_residual}{Analysis and synthesis residual.}{Across the tested dimensions and deformation values, the orthogonality error remains at floating-point scale.}

\FigureWide{14_opposing_q_casimir}{The opposing q-Casimir in one packet atlas.}{The tridiagonal form is the operator source of the q-Racah three-term recurrence.}

\FlowChapter{An exact q-controlled beam splitter}

Take
\[
\kappa_1=\kappa_2=\kappa_3=\frac12,
\qquad
N=1.
\]

\begin{theorem}[Closed two-channel q-Racah transform]
The left-to-right recoupling matrix can be chosen as
\[
\boxed{
\Rcal_q=
\begin{pmatrix}
c_q&s_q\\
s_q&-c_q
\end{pmatrix},
}
\]
where
\[
\boxed{
c_q=\frac{q}{1+q^2}}
\]
and
\[
\boxed{
s_q=\frac{\sqrt{1+q^2+q^4}}{1+q^2}.}
\]
\end{theorem}

\begin{proof}
The fixed-excitation space is spanned by $\ket{100},\ket{010},\ket{001}$. The total lowering equation is
\[
q x_1+x_2+q^{-1}x_3=0,
\]
so the lowest-weight multiplicity space is two-dimensional. Restricting either pair Casimir to this plane gives two normalized eigenvectors. Their overlap simplifies to the displayed matrix. Directly,
\[
c_q^2+s_q^2
=
\frac{q^2+1+q^2+q^4}{(1+q^2)^2}=1,
\]
so the matrix is orthogonal. Substitution into the restricted pair-Casimir eigenvalue equations completes the verification.
\end{proof}

At $q=1$,
\[
\Rcal_1=
\begin{pmatrix}
\frac12&\frac{\sqrt3}{2}\\[0.3em]
\frac{\sqrt3}{2}&-\frac12
\end{pmatrix},
\]
which is the classical equal-shape Racah beam splitter. As $q\downarrow0$,
\[
\Rcal_q\longrightarrow
\begin{pmatrix}0&1\\1&0\end{pmatrix}.
\]
The deformation therefore interpolates continuously between classical mixing and sparse routing.

\FigureWide{07_q_beam_splitter_amplitudes}{Closed beam-splitter amplitudes.}{The diagonal amplitude grows from zero to one half, while the transfer amplitude decreases from one to $\sqrt3/2$.}

\FigureWide{08_q_beam_splitter_kernel}{Squared two-channel kernels.}{The coherent transform remains orthogonal, while the measured channel becomes a q-dependent bistochastic matrix.}

\FlowChapter{Bistochastic measurement and entropy}

Define
\[
\boxed{
P_{jm}(q)=\abs{\Rcal_N(q)_{jm}}^2.
}
\]

\begin{theorem}[Bistochastic q-Racah channel]
For every finite recoupling matrix,
\[
P_{jm}\ge0,
\qquad
\sum_mP_{jm}=1,
\qquad
\sum_jP_{jm}=1.
\]
Consequently, for every probability row vector $p$,
\[
\boxed{
\entropy(pP)\ge\entropy(p),
}
\]
where $\entropy$ is Shannon entropy.
\end{theorem}

\begin{proof}
The row and column sums are the squared norms of rows and columns of the orthogonal matrix. Every doubly stochastic matrix is a convex combination of permutation matrices. Shannon entropy is concave and invariant under permutations, so it cannot decrease.
\end{proof}

\FigureWide{12_q_racah_probability_kernels}{Bistochastic q-Racah kernels.}{Every amplitude matrix produces a classical measurement channel with unit row and column sums.}

\begin{theorem}[Monotone one-excitation mixing]
For the equal-shape one-excitation channel and $0<q\le1$, the row entropy and the spectral mixing gap increase with q.
\end{theorem}

\begin{proof}
The smaller row probability is
\[
p(q)=c_q^2=\frac{q^2}{(1+q^2)^2}.
\]
It increases from zero to $1/4$. The binary entropy
\[
h(p)=-p\log p-(1-p)\log(1-p)
\]
is increasing for $0\le p\le1/2$. The two eigenvalues of the channel are $1$ and $2p-1$, so the spectral gap is
\[
\boxed{
1-\abs{2p-1}=2p=\frac{2q^2}{(1+q^2)^2},
}
\]
which is also increasing on $(0,1]$.
\end{proof}

\FigureWide{09_beam_splitter_entropy_gap}{Entropy and spectral gap in the exact two-channel model.}{Both increase monotonically with q.}

\begin{boundary}[No universal entropy monotonicity in q]
For higher excitation, each fixed q still gives a bistochastic channel and therefore increases entropy relative to its input. However, the row entropy need not be monotone as q varies. At $N=3$ and equal half-shapes, the second row entropy decreases from approximately $1.24488$ to $1.23710$ between $q\approx0.7772$ and $q\approx0.7884$. This negative control prevents the one-excitation theorem from being promoted to all packet levels.
\end{boundary}

\FigureWide{21_entropy_nonmonotonicity}{Higher-excitation entropy counterexample.}{A selected row of the q-Racah measurement channel is not monotone throughout the deformation interval.}

\FigureWide{22_stochastic_spectral_gap}{Spectral mixing beyond two channels.}{The q-dependence of the second eigenvalue becomes level dependent and need not follow a universal monotone law.}

\InlinePart{Deformation geometry}

\FlowChapter{Orthogonal transport in q}

Write
\[
q=\e^{-\theta},
\qquad
\theta\ge0.
\]
Choose a continuous sign gauge for the eigenvectors away from degeneracies and define
\[
R_N(\theta)=\Rcal_N(\e^{-\theta}).
\]

\begin{theorem}[Skew deformation generator]
The matrix
\[
\boxed{
A_N(\theta)
=R_N(\theta)^{\mathsf T}\partial_\theta R_N(\theta)
}
\]
is skew-symmetric, and
\[
\boxed{
\partial_\theta R_N=R_NA_N.
}
\]
\end{theorem}

\begin{proof}
Differentiate
\[
R_N^{\mathsf T}R_N=\Id.
\]
This gives
\[
(\partial_\theta R_N)^{\mathsf T}R_N
+R_N^{\mathsf T}\partial_\theta R_N=0,
\]
so $A_N^{\mathsf T}=-A_N$. Multiplying the definition of $A_N$ on the left by $R_N$ gives the tangent equation.
\end{proof}

Define the finite deformation metric
\[
\boxed{
g_N(\theta)=\frac12\norm{A_N(\theta)}_F^2.}
\]
It is invariant under a constant orthogonal change of either packet atlas. A q-dependent sign jump is not a smooth gauge and must be removed before differentiation.

\FigureWide{19_deformation_generator}{Skew generator of the q-flow.}{The antisymmetric matrix records which packet coordinates exchange amplitude under an infinitesimal change of deformation.}

\FigureWide{20_deformation_metric}{Deformation metric across excitation.}{The sensitivity of the complete recoupling atlas grows with multiplicity dimension.}

\FlowChapter{Classical and routing limits}

The q-Hahn and q-Racah systems have two useful endpoints.

At $q\to1$, symmetric q-numbers satisfy
\[
\qnum{x}=x+O((\log q)^2),
\]
so the packet algebra, Casimir spectra, lowest-weight vectors, and recoupling matrices converge to their ordinary $\mathfrak{su}(1,1)$ counterparts.

At $q\downarrow0$, no uniform finite limit exists for every unrenormalized q-number expression. Nevertheless, selected normalized packet transforms have finite routing limits. The one-excitation beam splitter is the simplest exact example.

\FigureWide{18_recoupling_classical_limit}{Classical limit of the matrix construction.}{After sign alignment, the complete q-Racah packet transform converges to the classical Racah transform.}

\InlinePart{Four-body q-associators}

\FlowChapter{Five complete q-Casimir tree atlases}

For four packets, the multiplicity dimension is
\[
\boxed{
\dim\Mcal_N^{(4)}=\binom{N+2}{2}.
}
\]
Consider the five binary coupling trees
\[
(((12)3)4),
\quad
((1(23))4),
\quad
(1((23)4)),
\quad
(1(2(34))),
\quad
((12)(34)).
\]
Each tree determines two commuting intermediate q-Casimirs. Their joint eigenvectors form a complete orthonormal basis $B_r$ of the same multiplicity space.

\FigureWide{23_q_pentagon_trees}{The q-deformed associahedral pentagon.}{Every vertex is a complete joint-Casimir atlas. Neighboring vertices are related by finite q-Racah packet transforms.}

\begin{theorem}[Finite q-pentagon certificate]
Let
\[
U_{r,r+1}=B_r^{\mathsf T}B_{r+1},
\]
with indices taken modulo five. Then
\[
\boxed{
U_{01}U_{12}U_{23}U_{34}U_{40}=\Id.
}
\]
\end{theorem}

\begin{proof}
Insert the definitions and use completeness:
\[
B_0^{\mathsf T}B_1
B_1^{\mathsf T}B_2
\cdots
B_4^{\mathsf T}B_0
=B_0^{\mathsf T}B_0=\Id.
\]
\end{proof}

This linear-algebra proof explains why the numerical pentagon is a strong implementation test. Any wrong projector, packet order, omitted vector, or inconsistent tree basis prevents the telescoping product from closing.

\FigureWide{24_q_pentagon_residual}{Pentagon residual over q and excitation.}{The closure error remains at floating-point scale through the tested multiplicity dimensions.}

\FigureWide{25_q_pentagon_edge_matrices}{Five q-associator edges.}{Each edge has its own signed amplitude pattern, but their ordered product reconstructs the initial tree atlas.}

\FlowChapter{Deforming the associator}

Each edge matrix depends on $\theta=-\log q$. In a continuous gauge, define
\[
A_r(\theta)
=U_r(\theta)^{\mathsf T}\partial_\theta U_r(\theta).
\]
Every $A_r$ is skew-symmetric. Differentiating the pentagon identity gives a transported sum of edge generators equal to zero. This is the finite flatness equation for the q-associator loop.

\FigureWide{26_associator_deformation}{Deformation of the five associator edges.}{The q-flow redistributes local generator norm around the pentagon while preserving exact global closure.}

\InlinePart{Finite compilation}

\FlowChapter{Givens factorization of a q-Racah packet}

\begin{theorem}[Finite rotation compiler]
Every real orthogonal $d\times d$ q-Racah packet matrix admits a factorization
\[
\boxed{
R=G_1^{\mathsf T}\cdots G_L^{\mathsf T}D,
}
\]
where $D$ is diagonal with entries $\pm1$ and every $G_j$ is a two-coordinate Givens rotation. One may choose
\[
\boxed{
L\le\frac{d(d-1)}2.
}
\]
\end{theorem}

\begin{proof}
Apply row Givens rotations from the lower left upward to eliminate all entries below the diagonal. Orthogonality forces the resulting upper triangular matrix to be diagonal with entries $\pm1$. Reversing the rotations reconstructs $R$.
\end{proof}

The compiler acts in the multiplicity packet. It does not require a dense gate acting on the full tensor-product level.

\FigureWide{27_q_givens_counts}{Rotation count.}{The direct compiler reaches the worst-case triangular count only when every subdiagonal elimination is needed.}

\FigureWide{28_q_givens_residual}{Compiler reconstruction residual.}{The finite rotation product reproduces the original q-Racah transform to numerical precision.}

\FigureWide{32_q_packet_compiler}{Programmable q-packet execution path.}{The compiler begins with shapes, q, and total excitation, constructs only the required multiplicity projectors, and emits signed transforms, measurement kernels, or pentagon certificates.}

\InlinePart{Scalar q-special functions}

\FlowChapter{Terminating q-Racah function}

For the q-Pochhammer symbol
\[
\qpoch{a}{n}=\prod_{j=0}^{n-1}(1-aq^j),
\]
define
\[
\boxed{
R_n(\mu(y);\alpha,\beta,\gamma,\delta\mid q)
={}_4\phi_3\!\left(
\begin{matrix}
q^{-n},\alpha\beta q^{n+1},q^{-y},\gamma\delta q^{y+1}\\
\alpha q,\beta\delta q,\gamma q
\end{matrix}
;q,q
\right),
}
\]
where
\[
\mu(y)=q^{-y}+\gamma\delta q^{y+1}.
\]
The series terminates at $n$. A finite orthogonal family arises when one denominator parameter is $q^{-N}$.

\begin{theorem}[q-Racah duality]
The terminating function satisfies
\[
\boxed{
R_n(\mu(y);\alpha,\beta,\gamma,\delta\mid q)
=
R_y(\mu(n);\gamma,\delta,\alpha,\beta\mid q).
}
\]
\end{theorem}

\begin{proof}
After exchanging $n$ with $y$ and $(\alpha,\beta)$ with $(\gamma,\delta)$, the four numerator q-Pochhammer factors are merely permuted. The three denominator factors are also permuted because $\delta\beta q=\beta\delta q$. Every summand is unchanged.
\end{proof}

\FigureWide{15_q_racah_polynomial_table}{A terminating q-Racah table.}{Degrees and lattice coordinates appear on equal footing, which makes duality visible before normalization by the orthogonality weights.}

\FigureWide{16_q_racah_duality}{Duality residual.}{The parameter-exchanged degree-coordinate identity holds at floating-point precision.}

\begin{theorem}[Termwise classical limit]
Set
\[
\alpha=q^A,
\quad
\beta=q^B,
\quad
\gamma=q^G,
\quad
\delta=q^D.
\]
For fixed $n$ and $y$,
\[
\boxed{
\lim_{q\to1^-}
R_n(\mu(y);q^A,q^B,q^G,q^D\mid q)
={}_4F_3\!\left(
\begin{matrix}
-n,n+A+B+1,-y,y+G+D+1\\
A+1,B+D+1,G+1
\end{matrix};1
\right).
}
\]
\end{theorem}

\begin{proof}
The q-series is finite. For every fixed integer $r$,
\[
\frac{\qpoch{q^x}{r}}{(1-q)^r}\longrightarrow(x)_r.
\]
Apply this limit to each q-Pochhammer factor in every summand and cancel the equal powers of $1-q$.
\end{proof}

\FigureWide{17_q_racah_to_racah}{q-Racah to Racah.}{The finite basic-hypergeometric polynomial converges termwise to its ordinary hypergeometric counterpart.}

\FlowChapter{Askey-Wilson transfer function}

For $z\in\C^\times$, define the normalized Askey-Wilson function
\[
\boxed{
\mathcal A_n(z)
={}_4\phi_3\!\left(
\begin{matrix}
q^{-n},abcdq^{n-1},az,az^{-1}\\
ab,ac,ad
\end{matrix};q,q
\right).
}
\]
Define the q-difference operator
\[
(\mathcal Lf)(z)
=A(z)(f(qz)-f(z))
+A(z^{-1})(f(q^{-1}z)-f(z))
+\left(1+\frac{abcd}{q}\right)f(z),
\]
where
\[
A(z)=
\frac{(1-az)(1-bz)(1-cz)(1-dz)}{(1-z^2)(1-qz^2)}.
\]

\begin{theorem}[Askey-Wilson eigenfunction law]
For every nonnegative integer $n$,
\[
\boxed{
\mathcal L\mathcal A_n
=
\left(q^{-n}+abcdq^{n-1}\right)\mathcal A_n.
}
\]
\end{theorem}

This is the continuous spectral lift of the finite q-Racah transfer system. The q-Racah family lives on a finite quadratic q-lattice. The Askey-Wilson family lives on the continuous variable $x=(z+z^{-1})/2$ and is governed by a second-order q-difference operator.

\FigureWide{29_askey_wilson_curves}{Askey-Wilson packet functions.}{Several degrees are plotted on the unit-circle coordinate $z=\e^{\ii\theta}$.}

\FigureWide{30_askey_wilson_difference_residual}{q-difference certificate.}{Direct substitution of the generated functions into the Askey-Wilson operator leaves only numerical roundoff.}

\FigureWide{31_q_function_class_map}{Finite and continuous q-functions.}{q-Hahn packets generate coupled bases, q-Racah functions generate finite recoupling, and Askey-Wilson functions provide the continuous spectral lift.}

\InlinePart{Quantum and computational interpretation}

\FlowChapter{A reversible q-shape module}

The q-Racah transform can be interpreted as a finite coherent basis change. The signed amplitude matrix is reversible. Squaring the entries and measuring the target basis produces the bistochastic channel.

The deformation parameter can be treated as a compile-time or control-time variable. A program may request:
\begin{enumerate}[leftmargin=2em]
\item a lowest-weight q-Hahn seed,
\item a left or right pair-Casimir atlas,
\item a signed q-Racah transform,
\item its measured bistochastic kernel,
\item a q-deformation tangent generator,
\item a four-body associator certificate,
\item a finite Givens circuit.
\end{enumerate}

This module can coexist with the logarithmic number Hamiltonian
\[
H\ket n=\log n\ket n
\]
used in the attached zeta-governed quantum dynamics. The logarithmic system supplies arithmetic time phases. The q-Racah system supplies finite shape coupling. Their tensor product is mathematically legitimate, but no theorem here identifies the q-Casimir spectrum with the nontrivial zeta zeros.

\FlowChapter{What the deformation adds}

The undeformed theory already has exact Racah recoupling and pentagon coherence. The q-coordinate adds four further layers.

First, it creates a continuous path between distinct packet-routing geometries.

Second, it turns basis sensitivity into the explicit skew generator $A_N(\theta)$.

Third, it produces q-dependent bistochastic measurement channels whose mixing behavior can be studied independently of coherent reversibility.

Fourth, it places the finite q-Racah packet and the continuous Askey-Wilson function inside one q-difference hierarchy.

The paper's strongest new operational object is therefore not a renamed polynomial. It is the complete tuple
\[
\boxed{
\left(
\Mcal_N,
\Cas_{12},
\Cas_{23},
\Rcal_N(q),
P_N(q),
A_N(\theta),
\{U_r(q)\}_{r=0}^{4}
\right).
}
\]
This tuple retains the object space, the two coordinate atlases, exact synthesis, the measured channel, the deformation tangent, and the associator loop.

\InlinePart{Verification and formulas}

\FlowChapter{Formula registry}

\begin{longtable}{L{0.29\textwidth}L{0.62\textwidth}}
\toprule
Object & Exact formula or certificate\\
\midrule
\endfirsthead
\toprule
Object & Exact formula or certificate\\
\midrule
\endhead
Symmetric q-number & $\qnum{x}=(q^x-q^{-x})/(q-q^{-1})$\\
Local q-algebra & $[K_-,K_+]=\qnum{2K_0}$\\
Coproduct & $\Delta K_\pm=K_\pm\otimes q^{K_0}+q^{-K_0}\otimes K_\pm$\\
q-Casimir & $\Cas_q=\qnum{K_0}\qnum{K_0-1}-K_+K_-$\\
Scalar Casimir & $\Cas_q=\qnum\kappa\qnum{\kappa-1}\Id$\\
Multiplicity dimension & $\dim\Mcal_N^{(m)}=\binom{N+m-2}{m-2}$\\
q-Hahn ratio & $\begin{aligned}
\frac{c_{n+1}}{c_n}={}&-q^{-(N-1+\kappa_1+\kappa_2)}\\
&\times\sqrt{\frac{\qnum{N-n}\qnum{2\kappa_2+N-n-1}}{\qnum{n+1}\qnum{2\kappa_1+n}}}
\end{aligned}$\\
Pair spectrum & $\lambda_j^{ab}=\qnum{\kappa_a+\kappa_b+j}\qnum{\kappa_a+\kappa_b+j-1}$\\
Recoupling & $\Rcal_N=L_N^{\mathsf T}R_N$\\
Bistochastic kernel & $P_{jm}=\abs{\Rcal_{jm}}^2$\\
Deformation generator & $A_N=\Rcal_N^{\mathsf T}\partial_\theta\Rcal_N$, $A_N^{\mathsf T}=-A_N$\\
Pentagon & $U_{01}U_{12}U_{23}U_{34}U_{40}=\Id$\\
q-Racah duality & $R_n(y;\alpha,\beta,\gamma,\delta)=R_y(n;\gamma,\delta,\alpha,\beta)$\\
Askey-Wilson law & $\mathcal L\mathcal A_n=(q^{-n}+abcdq^{n-1})\mathcal A_n$\\
Givens bound & $L\le d(d-1)/2$\\
\bottomrule
\end{longtable}

\FlowChapter{Numerical certificate table}

\begin{longtable}{L{0.42\textwidth}L{0.16\textwidth}L{0.28\textwidth}}
\toprule
Check & Result & Largest recorded residual or witness\\
\midrule
\endfirsthead
\toprule
Check & Result & Largest recorded residual or witness\\
\midrule
\endhead
Local q-algebra & pass & $2.05\times10^{-12}$\\
q-Hahn lowest vector & pass & $2.22\times10^{-16}$\\
Scalar q-Casimir & pass & $5.88\times10^{-15}$\\
q-Racah orthogonality & pass & $2.22\times10^{-15}$\\
Pair spectrum & pass & $4.55\times10^{-13}$\\
Opposing-Casimir tridiagonality & pass & $5.21\times10^{-12}$\\
Closed q-beam splitter & pass & $2.22\times10^{-16}$\\
Bistochastic kernel & pass & $2.00\times10^{-15}$\\
Entropy nondecrease for fixed kernel & pass & no violation in 80 trials\\
One-excitation entropy monotonicity & pass & no negative increment\\
Higher-level monotonicity control & pass & decrement $-7.78\times10^{-3}$ found\\
Matrix classical limit & pass & $2.82\times10^{-7}$\\
q-Hahn classical limit & pass & $4.28\times10^{-6}$\\
Routing limit & pass & $1.00\times10^{-5}$\\
Skew generator & pass & $1.48\times10^{-8}$ raw gauge residual\\
Tangent equation & pass & $5.92\times10^{-9}$\\
Four-tree commutation & pass & $7.96\times10^{-13}$\\
Joint diagonalization & pass & $7.33\times10^{-14}$\\
q-pentagon & pass & $3.22\times10^{-15}$\\
Givens compilation & pass & $8.88\times10^{-16}$\\
q-Racah duality & pass & $2.84\times10^{-14}$\\
q-Racah polynomial limit & pass & $4.27\times10^{-6}$\\
Askey-Wilson q-difference equation & pass & $5.91\times10^{-13}$\\
\bottomrule
\end{longtable}

The independent test suite repeats twenty properties through direct calls to the mathematical functions. It does not merely read the built-in check ledger.

\FlowChapter{Research directions and boundaries}

The present paper suggests several focused questions.

\begin{enumerate}[leftmargin=2em]
\item Determine closed shape-parameter formulas for the full q-Racah overlap in the symmetric q-number convention used here.
\item Characterize the deformation metric $g_N(\theta)$ asymptotically as $N$ grows.
\item Identify the precise packet conditions under which row entropy is monotone in q.
\item Extend the finite q-Casimir construction to continuous-series Askey-Wilson transforms.
\item Develop a q-adic or root-of-unity packet theory where semisimplicity may fail.
\item Compile the associator loop with locality constraints rather than unrestricted Givens rotations.
\item Couple the q-shape module to the logarithmic number Hamiltonian and determine which observables remain exactly factorized.
\end{enumerate}

The first five items are special-function and representation-theoretic questions. The last two are architecture and quantum-model questions. None is solved by the numerical figures alone.

\FlowChapter{Conclusion}

The q-deformation does not discard the earlier packet calculus. It enriches it.

The seed is still a positive shape packet. The object space is still resolved by a lowest-weight projector. Intermediate Casimirs still create exact coordinate atlases. The recoupling matrix still supplies analysis and synthesis. The squared entries still define a measurement channel. Four coupling trees still close around a pentagon. The finite transform still compiles into two-coordinate rotations.

What changes is the geometry of every layer. q alters the ladder, pair spectra, packet weights, recoupling amplitudes, stochastic mixing, and associator edges. The q-Racah and Askey-Wilson families are the scalar special functions that record this deformation.

The result is a complete q-packet module:
\[
\boxed{
\text{q-Hahn seed}
\longrightarrow
\text{q-Casimir atlases}
\longrightarrow
\text{q-Racah transform}
\longrightarrow
\text{q-pentagon}
\longrightarrow
\text{Askey-Wilson lift}.
}
\]
The classical theory is recovered at $q=1$. Selected packet transforms acquire sparse routing limits as $q$ approaches zero. Between those endpoints, the skew generator provides a local differential description of the changing coordinate system.

\appendix

\FlowChapter{Selected numerical tables}

\begin{table}[H]
\centering
\caption{Equal-shape one-excitation q-beam splitter.}
\begin{tabular}{ccccc}
\toprule
$q$ & $c_q$ & $s_q$ & row entropy & spectral gap\\
\midrule
0.10 & 0.099010 & 0.995086 & 0.055094 & 0.019606\\
0.25 & 0.235294 & 0.971924 & 0.214014 & 0.110727\\
0.50 & 0.400000 & 0.916515 & 0.439670 & 0.320000\\
0.75 & 0.480000 & 0.877268 & 0.539759 & 0.460800\\
1.00 & 0.500000 & 0.866025 & 0.562335 & 0.500000\\
\bottomrule
\end{tabular}
\end{table}

\begin{table}[H]
\centering
\caption{Four-body q-pentagon residuals at $N=4$.}
\begin{tabular}{ccc}
\toprule
$q$ & multiplicity dimension & maximum residual\\
\midrule
0.50 & 15 & $3.22\times10^{-15}$\\
0.70 & 15 & $4.00\times10^{-15}$\\
0.85 & 15 & $3.55\times10^{-15}$\\
1.00 & 15 & $3.44\times10^{-15}$\\
\bottomrule
\end{tabular}
\end{table}

\FlowChapter{Complete Python implementation}

The following source generates every figure, numerical table, built-in check, and manifest used in this paper.

\lstinputlisting{projector_generated_special_functions_viii.py}

\FlowChapter{Bibliographic note}

\begin{thebibliography}{9}

\bibitem{DLMFAskeyWilson}
NIST Digital Library of Mathematical Functions,
Section 18.28, ``Askey-Wilson Class.''

\bibitem{KoelinkVanDerJeugt}
H. T. Koelink and J. Van der Jeugt,
``Convolutions for orthogonal polynomials from Lie and quantum algebra representations,''
SIAM Journal on Mathematical Analysis 29 (1998), 794--822.

\bibitem{KoelinkRosengren}
E. Koelink and H. Rosengren,
``Harmonic analysis on the $SU(2)$ dynamical quantum group,''
Acta Applicandae Mathematicae 69 (2001), 163--220.

\bibitem{Shibukawa}
Y. Shibukawa,
``Clebsch--Gordan Coefficients for $\mathcal U_q(\mathfrak{su}(1,1))$ and $\mathcal U_q(\mathfrak{sl}(2))$, and Linearization Formula of Matrix Elements,''
Publ. Res. Inst. Math. Sci. 28 (1992), 775--807.

\bibitem{GenestIlievVinet}
V. X. Genest, P. Iliev, and L. Vinet,
``Coupling coefficients of $su_q(1,1)$ and multivariate q-Racah polynomials,''
Nuclear Physics B 927 (2018), 97--123.

\bibitem{GasperRahman}
G. Gasper and M. Rahman,
\emph{Basic Hypergeometric Series}, second edition,
Cambridge University Press, 2004.

\bibitem{WeiEtAl}
S. Wei, Y. Zhai, Q. Lu, et al.,
``The Riemann Hypothesis manifested in dynamical quantum phase transitions,''
Nature Communications, accepted 18 June 2026,
DOI 10.1038/s41467-026-74935-8.

\end{thebibliography}

\FlowChapter{Reproduction}

From the source directory, run
\begin{lstlisting}[language=bash,numbers=none]
python projector_generated_special_functions_viii.py
python -m pytest -q
pdflatex -interaction=nonstopmode Projector_Generated_Special_Functions_VIII.tex
pdflatex -interaction=nonstopmode Projector_Generated_Special_Functions_VIII.tex
pdflatex -interaction=nonstopmode Projector_Generated_Special_Functions_VIII.tex
\end{lstlisting}

The Python step writes all 32 figure pairs, the check ledger, the selected data tables, and the machine-readable manifest before the paper is compiled.

\end{document}
